\section{Reversible changes of the coagulation kernel}
\label{sec:kernel-comparison}

The stationary law can be preserved while allowing the reaction rates
to depend on cluster sizes.  Fix positive symmetric multipliers
$c_n(i,j)=c_n(j,i)$ for $i+j\le n$.
An unordered pair of cluster instances of sizes $i,j$ merges at
rate $2c_n(i,j)/n$.  Each internal link whose removal produces lengths
$i,j$ fragments at rate $\beta c_n(i,j)$.
Thus, for $i<j$, the coarse forward and reverse rates are
\[
 \frac{2c_n(i,j)}n m_im_j,
 \qquad 2\beta c_n(i,j)m_{i+j},
\]
whereas for $i=j$ they are
\[
 \frac{c_n(i,i)}n m_i(m_i-1),
 \qquad \beta c_n(i,i)m_{2i}.
\]

\begin{theorem}\label{thm:kernel-comparison}
The modified chain is reversible for the original law $\pi_M$ in
\eqref{proof:eq:piM}.  If
\[
 0<c_-\le c_n(i,j)\le c_+<\infty
 \qquad(i+j\le n),
\]
then its gap $\gamma_c$ satisfies
\begin{equation}\label{eq:kernel-gap-comparison}
 c_-\gamma_M(n,\beta)\le\gamma_c(n,\beta)
      \le c_+\gamma_M(n,\beta).
\end{equation}
In particular, with $s=\beta+\sqrt\beta+n^{-1}$,
\[
 \frac{c_-}{6}s\le\gamma_c\le4c_+s.
\]
The comparison constants are uniform in $n,\beta$ whenever
$c_-,c_+$ are.
\end{theorem}

\begin{proof}
Each reaction and its reverse have been multiplied by the same
positive number.  Every original detailed-balance identity therefore
remains valid.  Writing the Dirichlet form as a sum over undirected
reaction edges gives
\[
 c_-\mathcal E_M(f,f)\le\mathcal E_c(f,f)
       \le c_+\mathcal E_M(f,f).
\]
The variance uses the same stationary law for all three forms.
Taking infima of the Rayleigh quotients proves
\eqref{eq:kernel-gap-comparison}.
\end{proof}

This is a robustness consequence of Dirichlet-form comparison.
Its proof also applies to symmetric multipliers attached to full
undirected state-space edges.
It does not establish the sharp scale of an unbounded kernel from
its typical value under equilibrium.

The stationary cut representation survives without a boundedness
assumption on $c_n$.  At a boundary between segments of lengths $i,j$,
replace the deletion rate by $(k-1)c_n(i,j)/n$.
At an internal link producing $i,j$, use the birth rate
$\beta c_n(i,j)$.  The law $\pi_B$ is unchanged, and averaging over
the uniform cyclic order conditional on the segment multiset gives
the modified polymer generator.  Consequently the exact variance
and Dirichlet identities in \eqref{proof:eq:dirichletcomparison}
continue to hold.

The full harmonic reduction generally does not survive.
For example, take $n=4$ and $c(i,j)=i+j$.
Cut configurations with segment lengths $(1,3)$ and $(2,2)$ both have
$k=2$, but their generator applied to $k$ is, respectively,
$6\beta-2$ and $4\beta-2$.
Even the subspace of functions of cardinality is therefore not
invariant for this modified cut generator.

\subsection{A classical multiplicative-kernel comparison class}

A different Gibbs equilibrium gives a useful growing-kernel baseline.
Consider pair-merger rates $2ij/n$ and per-link fragmentation rate
$\beta(i+j)$ when the two resulting lengths are $i,j$.
Writing $\theta=\beta n$, the reversible law is
\begin{equation}\label{eq:ewens-polymer-law}
 \pi_{\mathrm E}(m)=Z_{\mathrm E}^{-1}
    \frac{\theta^{K(m)}}{\prod_{j=1}^n j^{m_j}m_j!}.
\end{equation}
This is the Ewens cycle-count law.
The chain is the classical Metropolis transposition process on
cycle types, with a deterministic change of time.

\begin{proposition}[Classical spectral formula in the present units]
\label{prop:ewens-classical}
For this multiplicative-kernel chain,
\begin{equation}\label{eq:ewens-gap}
 \gamma_{\mathrm E}(n,\beta)=2\beta+\frac{2(n-1)}n.
\end{equation}
Multiplying both directions of each size-pair reaction by
$c_n(i,j)\in[c_-,c_+]$ preserves \eqref{eq:ewens-polymer-law} and gives
\[
 c_-\left(2\beta+\frac{2(n-1)}n\right)
 \le\gamma_{\mathrm E,c}\le
 c_+\left(2\beta+\frac{2(n-1)}n\right).
\]
\end{proposition}

\begin{proof}
Detailed balance in \eqref{eq:ewens-polymer-law} follows directly from
the weight ratio $w_iw_j/w_{i+j}=(i+j)/(ij)$ for $w_j=1/j$.
For unequal fragments the coarse split rate is
$2\beta(i+j)m_{i+j}$, and for equal fragments it is
$\beta(2i)m_{2i}$.

We record the normalization of the known spectrum explicitly.
In Jiang~\cite[Eq.~(7) and Theorem~4.2]{Jiang2015}, the parameter
of the Metropolis transition matrix is $t=\theta^{-1}$.
Our continuous-time generator is
$(n-1)\max\{1,\theta\}(P_t-I)$.
Indeed, multiplying the uniform transposition proposal rate
$\binom n2^{-1}$ by this factor gives merger rate $2/n$ per pair
of labels in distinct cycles and split rate $2\theta/n$ per pair
of labels in one cycle.
Jiang's spectral formula consequently gives, for each partition
$\lambda\vdash n$,
\[
 \Lambda_\lambda
 =\frac2n\left[
      \binom n2-\sum_i\binom{\lambda_i}{2}
       +\theta\sum_i(i-1)\lambda_i\right].
\]
The zero eigenvalue corresponds to $\lambda=(n)$.
For every other partition,
$\binom n2-\sum_i\binom{\lambda_i}{2}\ge n-1$ and
$\sum_i(i-1)\lambda_i\ge1$.
Both are equalities for $\lambda=(n-1,1)$.
This proves \eqref{eq:ewens-gap}; the comparison statement follows
by the preceding Dirichlet-form argument.
\end{proof}

The spectral formula in Proposition~\ref{prop:ewens-classical}
is an application of an established diagonalization, not a new
solution of the multiplicative-kernel model.
Its perturbation class illustrates that rates growing like $ij$
can have a uniform gap of order $1+\beta$ once the fragmentation
rule and equilibrium law are specified together.
